\section{Introducción}
El café es un producto de alcance mundial: la Organización Internacional del Café \citeyearpar{ico2026} estimó una producción de 177,5 millones de sacos en el año cafetero 2024 y un consumo de 175,1 millones, y el Servicio Agrícola Exterior del Departamento de Agricultura de Estados Unidos \citeyearpar{usdafas2026} proyecta para 2026/27 una producción récord de 189,7 millones de sacos y un consumo de 179,7 millones. En Bolivia, en 2025 el valor exportado fue récord, con más de 17 millones de dólares y un aumento de 13\,\% frente a 2024, aunque el volumen bajó 12\,\% y 99\,\% de lo exportado fue café sin tostar \citep{ibce2026}.

En ese contexto, el nivel de tueste define el color, el sabor y el valor del producto. La norma del Servicio de Comercialización Agrícola del Departamento de Agricultura de Estados Unidos \citeyearpar{usdaams2004} distingue cuatro tonos medidos con la escala Agtron: claro (80 a 90), medio (50 a 60), moderadamente oscuro (40 a 50) y oscuro (30 a 40), y pide que el color de los granos coincida con las lecturas del colorímetro. El tueste también pesa en el ingreso: en los Yungas, \citet{jacobi2024} registraron que el café especial verde se pagó a 2,69 dólares por libra, mientras que el tostado alcanzó entre 10,96 y 11,84 dólares por libra según el canal de venta.

Evaluar el tueste a simple vista depende de la persona: \citet{marzuki2025} señalan que la evaluación humana se ve afectada por la fatiga, la iluminación y la experiencia personal, y proponen una alternativa objetiva con imágenes. La visión artificial ya se aplicó a esta tarea: Marzuki et al. lograron 96,75\,\% de exactitud de validación con una red convolucional y cuatro niveles de tueste; \citet{anto2025} compararon cinco arquitecturas con ocho niveles Agtron sobre 2.522 imágenes tomadas con luz controlada y obtuvieron hasta 93,8\,\%; y \citet{arsyan2026} alcanzaron 96\,\% con MobileNetV2 preentrenada y cuatro niveles.

Estas redes compactas se pueden reutilizar mediante aprendizaje por transferencia \citep{howard2017,sandler2018,yosinski2014}, que es la base de herramientas sin código como Teachable Machine. Sin embargo, una exactitud alta en la prueba no siempre implica que el modelo aprendió lo que se pretendía. \citet{geirhos2020} describen los atajos como reglas de decisión que rinden bien en las pruebas habituales pero fallan en condiciones más exigentes, como las del uso real; \citet{xiao2020} muestran que el fondo de la imagen basta para clasificar con éxito razonable, y \citet{torralba2011} reportan que un clasificador de automóviles con 53,4\,\% de exactitud en su propio conjunto de datos baja a 27,5\,\% al probarlo con otros. A ello se suma la fuga de datos, que \citet{kapoor2022} identificaron en 329 artículos de 17 campos. En los estudios sobre café, \citet{arsyan2026} proponen ampliar el conjunto de imágenes con más fondos y condiciones de luz.

Este manuscrito resulta de la primera etapa de un proyecto de la asignatura Inteligencia Artificial (ISW-221) de la Universidad Católica Boliviana. El alcance de esta etapa se limita a tomar fotografías de los granos, cargarlas en Teachable Machine, entrenar el modelo y entregar un prototipo funcional; la acción física con un microcontrolador ESP32 y los requisitos que defina la asignatura corresponden a la etapa siguiente. El objetivo fue construir un prototipo que clasifique el nivel de tueste del café a partir de imágenes con Teachable Machine y comprobar si el modelo se apoya en el grano o en el fondo.
